\documentclass{jmlr}
% jmlr.cls auto-loads: amsmath, amssymb, natbib, graphicx, url
\usepackage{float}
\usepackage{amsthm,bm}
\usepackage{nameref,hyperref}
\usepackage{lastpage}
\usepackage{epstopdf}
\usepackage{booktabs}
\usepackage{subcaption}
\usepackage{longtable}
\usepackage{array}
\usepackage[T1]{fontenc}
\usepackage{lmodern}
\usepackage{textcomp}
\usepackage{pgfplots}
\pgfplotsset{compat=1.18}
\usepackage{tikz}
\usetikzlibrary{arrows.meta, decorations.pathreplating}

% ---- JMLR metadata (fill in after acceptance) --------------------------------
\jmlrvolume{vol}
\jmlryear{year}
\jmlrsubmitted{submitted}
\jmlrpublished{published}

\title{Understanding When Poisson Log-Normal Models Outperform Penalized Poisson Regression for Microbiome Count Data}

\author{%
  \Name{Daniel Agyapong} \Email{da2343@nau.edu}\\
  \addr School of Informatics, Computing, and Cyber Systems,
  Northern Arizona University, Flagstaff, Arizona, USA
  \AND
  \Name{Julien Chiquet} \Email{julien.chiquet@inrae.fr}\\
  \addr UMR MIA Paris-Saclay, Universit\'{e} Paris-Saclay,
  AgroParisTech, INRAE, Palaiseau, France
  \AND
  \Name{Jane Marks} \Email{jane.marks@nau.edu}\\
  \addr Department of Biological Sciences,
  Northern Arizona University, Flagstaff, Arizona, USA
  \AND
  \Name{Toby Dylan Hocking} \Email{toby.hocking@usherbrooke.ca}\\
  \addr D\'{e}partement d'informatique, Universit\'{e} de Sherbrooke,
  Sherbrooke, Qu\'{e}bec, Canada%
}

\editor{Editor}

\begin{document}
\maketitle

\begin{abstract}
Poisson Log-Normal (PLN) models are increasingly applied to microbiome count data for prediction and network inference, yet systematic evidence on when they outperform simpler and faster penalized Poisson regression remains scarce. 
The PLN family introduces a latent Gaussian layer to capture extra Poisson variability, supporting conditional prediction, clustering, and network inference, but their computational demands and sample size requirements raise the practical question of when this added complexity is warranted. 
We address this through a systematic empirical benchmark comparing PLN and \texttt{GLMNet(Poisson)} across 34 real microbiome datasets spanning sample sizes $N = 32$--$18{,}270$ and dimensionalities $D = 20$--$373$.
Prediction is evaluated through a Leave-One-Taxon-Out cross-validation scheme using held-out Poisson deviance, while network inference is evaluated against experimentally validated interaction datasets using F1 score.
For prediction, PLN outperforms \texttt{GLMNet(Poisson)} on the majority of datasets, with improvements reaching up to 38\%, driven primarily by the $N/D$ ratio and count overdispersion.
PLN wins when $N/D$ is low and counts are substantially overdispersed, where the latent Gaussian layer captures genuine extra-Poisson variance; \texttt{GLMNet(Poisson)} dominates when $N/D$ is large or overdispersion is low.
For network inference, PLNNetwork consistently outperforms \texttt{GLMNet} neighbourhood selection on broad, undirected interaction networks, while \texttt{GLMNet(Poisson)} is better suited to local or directional effects.
Together, these results provide the first unified empirical map of when PLN's latent dependence modeling confers a genuine advantage over penalized Poisson regression.

\textbf{Keywords:} microbiome, Poisson Log-Normal, penalized regression, network inference,\\ cross-validation, benchmarking 

\end{abstract}

\section{Introduction}
Microbial communities profiled by Next-Generation Sequencing produce count data with distinctive statistical properties.
Observations are nonnegative integers, contain many zeros, and are strongly overdispersed.   
Moreover, the number of taxa $D$ routinely exceeds the number of biological samples $N$~\cite{lagkouvardos2016imngs, Kodikara2022}.
These properties jointly challenge the assumptions of standard Gaussian models and limit the use of conventional Poisson regression, which cannot accommodate residual correlations among taxa or the extra-Poisson variability that arises from unobserved biological heterogeneity~\cite{mcmurdie2014waste, Chiquet2021}.

The Poisson Log-Normal (PLN) model addresses these limitations by introducing a latent Gaussian layer that captures both overdispersion and dependence among taxa within a principled probabilistic framework~\cite{aitchison1989multivariate, Chiquet2021}.
This latent structure supports a family of downstream analyses, including conditional prediction, dimensionality reduction, and the recovery of sparse ecological interaction networks through graphical lasso regularization~\cite{chiquet2019variational}.
The PLN family is fitted using variational inference and has been increasingly applied in microbiome studies that require explicit modeling of taxon dependencies~\cite{Chiquet2021, Subedi2025}.

Despite this appeal, PLN models carry a practical cost.
Estimating a full $D \times D$ latent covariance matrix scales cubically in $D$, making PLN computationally demanding for the large, sparse count tables typical of modern microbiome surveys.
By contrast, penalized Poisson regression implemented efficiently in \texttt{glmnet}~\cite{Friedman2010} fits each taxon independently using $\ell_1$ regularization, scales near-linearly in $D$, and remains tractable even when $D$ runs into the thousands.
The fundamental question motivating this work is therefore empirical.
Under what conditions does PLN's added modeling complexity translate into measurable gains over penalized Poisson regression?

This question has direct practical consequences.
A researcher analyzing a dataset with $N = 50$ samples and $D = 200$ taxa faces a genuine modeling choice between investing in PLN's joint estimation at computational cost and relying on the faster penalized regression baseline.
Existing literature provides limited guidance.
Prior evaluations of PLN variants have largely focused on illustrative case studies, synthetic data, or narrow analytical tasks, assessing goodness of fit on training data rather than held-out predictive performance~\cite{Chiquet2021, Batardiere2025ZIPLN, Chaussard2025PLNTree}.
As summarized in Table~\ref{tab:related_work}, systematic comparisons against penalized regression baselines on real microbiome data under a unified and reproducible evaluation protocol remain absent from the literature.

A similar gap exists for network inference.
PLNNetwork recovers ecological interaction networks by estimating a sparse latent precision matrix~\cite{chiquet2019variational}, while neighbourhood selection through \texttt{GLMNet(Poisson)} offers an alternative route to network recovery without a joint model~\cite{Meinshausen2006, Friedman2010}.
To our knowledge, no study has evaluated these competing approaches against experimentally validated microbial interactions, which are necessary for a principled comparison of edge recovery performance.

In this paper, we present a systematic empirical benchmark addressing both tasks.
For prediction, we compare PLN and \texttt{GLMNet(Poisson)} across 34 real microbiome datasets under a Leave-One-Taxon-Out five-fold cross-validation (LOTO-CV) scheme~\cite{agyapong2025cross}, evaluating held-out Poisson deviance across a wide range of sample sizes and dimensionalities.
For network inference, we evaluate PLN with graphical lasso against \texttt{GLMNet} neighbourhood selection using five experimentally validated interaction datasets spanning defined microbial communities and in vivo systems.
Together, these experiments provide the first unified empirical map of when PLN's latent dependence modeling confers a genuine advantage over penalized Poisson regression.

The remainder of the paper is organized as follows.
Section~2 reviews related work and situates this study within the existing benchmark literature.
Section~3 describes the algorithms, datasets, and evaluation protocol.
Section~4 presents results for prediction and network inference.
Section~5 discusses computational scaling.
Section~6 concludes.

\section{Related Work}

Microbiome abundance data poses statistical challenges that motivate specialized count models; this section reviews the modelling choices that underlie PLN and the benchmarks that have evaluated it.
Table~\ref{tab:related_work} summarizes how prior studies approach evaluation.

\paragraph{Count data properties of microbiome surveys.}
Microbiome surveys via 16S rRNA amplicon sequencing or shotgun metagenomics produce read-count tables of non-negative integers whose observed values depend on the total number of reads sequenced per sample~\cite{HMP2012, Thompson2017, Xia2023}.
Three properties of these counts make standard Gaussian regression inappropriate. 
The latter two also show that plain Poisson regression is insufficient, motivating latent-variable extensions such as the Poisson Log-Normal.
First, \emph{compositionality}: total read depth (library size) varies across samples by orders of magnitude and carries no biological information; models that ignore this offset confound technical sequencing variation with biological signal~\cite{mcmurdie2014waste, Xia2023}.
Second, \emph{overdispersion}: counts exhibit variance far exceeding the Poisson mean, as genuine biological variation across individuals and environments inflates count variance beyond what the Poisson distribution allows~\cite{zhang2017nbmm, McCullaghNelder1989}. 
Third, \emph{excess zeros}: the majority of entries in the count matrix are zero not because taxa are truly absent but because rare taxa fall below the sequencing detection threshold, with data sparsity commonly exceeding 70\%~\cite{Batardiere2025ZIPLN, Kodikara2022}.
Gaussian regression treats counts as unbounded continuous quantities.
Ad hoc transformations such as $\log(1+x)$ can improve numerical behavior and stability, but they do not by themselves resolve discrete support or compositional structure~\cite{booeshaghi2021log1p}.

\paragraph{The Poisson Log-Normal model.}
The generalized linear model with Poisson or Negative Binomial likelihood is the natural framework for overdispersed count data~\cite{McCullaghNelder1989}.
For univariate differential abundance testing, Negative Binomial models are widely used, as implemented in DESeq2~\cite{Love2014DESeq2} and edgeR~\cite{Robinson2010edgeR}.
These are, however, marginal models: each taxon is fitted independently, so the joint dependence structure across taxa is not represented.
The Poisson Log-Normal (PLN) model, first introduced by Aitchison and Ho~\cite{aitchison1989multivariate}, addresses this by placing a multivariate Gaussian distribution on the log-intensities, capturing correlations between different taxa through a shared latent covariance matrix while still respecting the discrete, non-negative nature of the data and accommodating library-size offsets.
This makes PLN a single coherent framework for prediction, ordination, clustering, and network inference from count data~\cite{Chiquet2021, Subedi2025}.

\paragraph{Prediction and representation learning.}
Poisson PCA has been assessed primarily through simulations emphasizing principal component recovery, robustness, and computational efficiency, with microbiome datasets serving as motivating illustrations~\cite{Kenney2021PoissonPCA, Virta2023PoissonPCA}.
These studies evaluate reconstruction error and latent dimension recovery on training data rather than held-out predictive performance, and do not include penalized Poisson regression as a baseline.
The PLN framework paper~\cite{Chiquet2021} presents a suite of qualitative applications spanning ordination, regression, network inference, and clustering across several ecological and microbiome datasets, but does not impose a standardized cross-validation protocol or compare against penalized regression.
Extensions addressing zero inflation and hierarchical structure compare PLN variants on synthetic and real microbiome data using in-sample criteria such as variational evidence lower bound (ELBO), log-likelihood or information criteria (AIC/BIC), and downstream discrimination or generative performance~\cite{Batardiere2025ZIPLN, Chaussard2025PLNTree}.

\paragraph{Network inference.}
Network benchmarks for count data graphical models evaluate edge recovery on synthetic PLN graphs and single-cell or gene expression data, reporting ROC curves, AUPR, and edge-weight correlations~\cite{Xiao2022PLNet}.
Synthetic data generators drawing from the PLN distribution have been developed to stress-test network algorithms under controlled graph topologies and sparsity levels~\cite{Qian2024}.
More broadly, Ghaeli et al.~\cite{Ghaeli2025} proposed a benchmarking framework focused on reproducibility, using bootstrap resampling to generate modified versions of real microbiome data, but stopped short of validating inferred edges against experimentally confirmed interactions.
Across these studies, edge-recovery assessment relies on synthetic or statistically constructed ground truth.

\begin{table}[t]
\caption{Summary of prior benchmarks of PLN and Poisson algorithms.}
\centering
\small
\begin{tabular}{>{\raggedright\arraybackslash}p{2.8cm}>{\raggedright\arraybackslash}p{3.6cm}>{\raggedright\arraybackslash}p{3.6cm}>{\raggedright\arraybackslash}p{4.6cm}}
\toprule
Study (year) & Data / setting & Algorithms tested & Evaluation metrics \\
\midrule
Kenney et al. (2021)~\cite{Kenney2021PoissonPCA} & Simulations; microbiome example & Poisson PCA, PLN & Subspace recovery (principal angles); robustness to noise and exposure; runtime \\
\midrule
Chiquet et al. (2021)~\cite{Chiquet2021} & Multiple ecology/microbiome case studies & PLN, PLNPCA, PLNNetwork, PLNMixture & BIC/ICL for model order selection; illustrative fits across tasks; no held-out evaluation \\
\midrule
Xiao et al. (2022)~\cite{Xiao2022PLNet} & Synthetic count networks; real scRNA-seq data & PLNet, VPLN, latentcor, glasso & AUPR/AUC for edge recovery; runtime \\
\midrule
Virta \& Artemiou (2023)~\cite{Virta2023PoissonPCA} & Simulations; real count matrices & Poisson PCA, vectorization PCA, matrix-normal methods & Subspace estimation error (principal angles); latent-rank selection accuracy; runtime \\
\midrule
Qian et al. (2024)~\cite{Qian2024} & Synthetic microbial communities (5 sizes, 4 topologies) & PLN, MIC, Pearson, Ma et al. & 3-class accuracy for interaction type recovery; topology and size effects \\
\midrule
Batardi\`ere et al. (2025)~\cite{Batardiere2025ZIPLN} & High-zero simulations; cow microbiome & ZIPLN, PLN & In-sample log-likelihood (ELBO/BIC/AIC); latent dispersion; group discrimination \\
\midrule
Chaussard et al. (2025)~\cite{Chaussard2025PLNTree} & Synthetic datasets; human gut microbiome (T2D) & PLN-Tree, mean-field PLN, PLN & ELBO/log-likelihood; Shannon, Simpson, Bray-Curtis diversity; downstream classification accuracy \\
\midrule
\textbf{This work (2026)} & 34 real microbiome datasets; 5 experimental interaction datasets & PLN, PLNNetwork, \texttt{GLMNet(Poisson)}, featureless & Prediction: held-out Poisson deviance under 5-fold LOTO-CV; network inference: F1 against experimental ground truth; runtime; memory \\
\bottomrule
\end{tabular}
\label{tab:related_work}
\end{table}


\paragraph{Summary of contributions.}
Existing benchmarks of Poisson regression and Poisson Log Normal algorithms have primarily targeted narrow analytical tasks, evaluating individual variants such as Poisson PCA~\cite{Kenney2021PoissonPCA, Virta2023PoissonPCA}, PLN regression or network formulations~\cite{Chiquet2021, Xiao2022PLNet}, and zero-inflated or hierarchical extensions~\cite{Batardiere2025ZIPLN, Chaussard2025PLNTree} but typically in isolation, on synthetic data, or through qualitative case studies rather than unified systematic comparisons.
Most prior evaluations emphasize synthetic or illustrative demonstrations, lack standardized cross-validation protocols, omit uncertainty quantification, and rarely compare PLN variants directly to penalized Poisson regression algorithms on real microbiome data. 
As shown in Table~\ref{tab:related_work}, the field still lacks a unified empirical framework capable of assessing both predictive accuracy and model stability under consistent experimental conditions.

\begin{itemize}
  \item We benchmark PLN against \texttt{GLMNet(Poisson)} for count prediction across 34 real microbiome datasets under a unified Leave-One-Taxon-Out cross-validation~\cite{agyapong2025cross} protocol, quantifying predictive uncertainty through paired deviance comparisons. This provides the first held-out comparison of these two model classes on real data.
  \item We identify $N/D$ as the primary empirical predictor of relative performance, and overdispersion as the strongest secondary predictor. PLN improves on \texttt{GLMNet(Poisson)} by up to 38\% when $N/D$ is low and counts are strongly overdispersed; \texttt{GLMNet(Poisson)} dominates when $N \gg D$.
  \item We evaluate PLNNetwork and \texttt{GLMNet} neighbourhood selection against experimentally validated microbial interaction datasets, providing the first comparison of these approaches on real biological ground truth rather than simulated or statistically altered networks.
\end{itemize}

%%%%%%%%%%%%%%%%%%
\begin{figure}[t]
\centering
\begin{tikzpicture}[
  font=\small,
  pbox/.style={rectangle, draw=gray!55, fill=gray!10, rounded corners=3pt,
               align=center, minimum width=3.8cm, minimum height=0.88cm, font=\footnotesize},
  tbox/.style={rectangle, draw=blue!55, fill=blue!10, rounded corners=3pt,
               align=center, minimum width=3.8cm, minimum height=0.88cm, font=\footnotesize},
  parrow/.style={-Stealth, thick, gray!55},
  tarrow/.style={-Stealth, thick, blue!55},
]
\node[font=\small\bfseries, text=gray!70]       at (2.3, 8.5) {Prior work};
\node[font=\small\bfseries, text=blue!70!black] at (8.9, 8.5) {This work};
\draw[dashed, thick, gray!35] (5.6, -2.1) -- (5.6, 8.3);
\node[font=\footnotesize\itshape, text=gray!60]       at (2.3, 7.95) {prediction};
\node[font=\footnotesize\itshape, text=blue!55!black] at (8.9, 7.95) {prediction};
	\node[pbox] (p1) at (2.3, 7.3)  {Synthetic or real data\\(no holdout)};
	\node[pbox] (p2) at (2.3, 5.65) {PLN family};
	\node[pbox] (p3) at (2.3, 4.0)  {ELBO / BIC / AIC\\(in-sample)};
	\draw[parrow] (p1) -- (p2);
	\draw[parrow] (p2) -- (p3);
	\node[tbox] (t1) at (8.9, 7.3)  {28 real microbiome\\datasets};
	\node[tbox] (t2) at (8.9, 5.65) {PLN vs \texttt{GLMNet(Poisson)}};
	\node[tbox] (t3) at (8.9, 4.0)  {Held-out Poisson\\deviance (LOTO-CV)};
\draw[tarrow] (t1) -- (t2);
\draw[tarrow] (t2) -- (t3);
\draw[dashed, gray!28] (0.0, 3.2) -- (11.2, 3.2);
\node[font=\footnotesize\itshape, text=gray!60]       at (2.3, 2.85) {network inference};
\node[font=\footnotesize\itshape, text=blue!55!black] at (8.9, 2.85) {network inference};
	\node[pbox] (n1) at (2.3, 2.25)  {Simulated or\\modified networks};
	\node[pbox] (n2) at (2.3, 0.6)   {PLNNetwork};
	\node[pbox] (n3) at (2.3, -1.05) {ROC / AUPR on\\simulated ground truth};
	\draw[parrow] (n1) -- (n2);
	\draw[parrow] (n2) -- (n3);
	\node[tbox] (m1) at (8.9, 2.25)  {5 experimentally\\validated datasets};
\node[tbox] (m2) at (8.9, 0.6)   {PLNNetwork vs \texttt{GLMNet(Poisson)}};
\node[tbox] (m3) at (8.9, -1.05) {F1 vs validated\\microbial interactions};
\draw[tarrow] (m1) -- (m2);
\draw[tarrow] (m2) -- (m3);
\end{tikzpicture}
\caption{Evaluation pipelines for prediction (top) and network inference (bottom).
Prior work for prediction includes Chiquet et al.~\cite{Chiquet2021}, Batardi\`{e}re et al.~\cite{Batardiere2025ZIPLN}, and Chaussard et al.~\cite{Chaussard2025PLNTree}.
Prior work for network inference includes Xiao et al.~\cite{Xiao2022PLNet}, Qian et al.~\cite{Qian2024}, and Ghaeli et al.~\cite{Ghaeli2025}.}
\label{fig:pipeline}
\end{figure}
% \begin{figure}[t]
%     \centering
%     \includegraphics[width=\textwidth]{concept_nov.png}
%     \caption{Previous vs Proposed Evaluation Criteria.}
%     \label{fig:concept}
% \end{figure}

\section{Methodology}
\subsection{Real Microbiome Abundance datasets}
\label{subsec:count-prediction-datasets}
We evaluated count prediction on a heterogeneous benchmark of 28 microbiome datasets spanning oral, skin, nasal, stool, infant, pediatric, and general gut systems, together with several collections focused on disease including colorectal  (CRC), inflammatory bowel disease, \textit{Clostridium difficile} infection, HIV, non-alcoholic steatohepatitis, obesity, Parkinson's disease, rheumatoid arthritis, and type~1 diabetes.
The retained benchmark set includes feature tables at the family, genus, species, and OTU levels and covers a broad range of sample sizes and feature dimensions.
Across datasets, the number of samples ranged from $N=32$ to $N=18{,}364$, the number of features ranged from $D=24$ to $D=373$, and the fraction of zero entries in the benchmarked abundance matrices ranged from $30.64\%$ to $93.01\%$.

The benchmark collection is summarized in Table~\ref{tab:bmr-count-datasets}.
We retained the non-redundant dataset variants actually used in the count-prediction analysis and omitted near-duplicate and very small pilot datasets.
Full benchmark identifiers and provenance are preserved in the repository metadata for exact reproducibility.

\begin{table}[t]
\caption{Count-prediction datasets included in the BMR benchmark, shown with shortened manuscript-friendly labels.
Data sparsity is the percentage of zero entries in the benchmarked abundance matrix.
Citation markers refer to the corresponding primary study records or dataset portals.
Full dataset identifiers are retained in the repository metadata.}
\label{tab:bmr-count-datasets}
\centering
\small
\begin{tabular}{>{\raggedright\arraybackslash}p{0.24\textwidth}>{\raggedright\arraybackslash}p{0.22\textwidth}>{\raggedright\arraybackslash}p{0.08\textwidth}rrr}
\toprule
Dataset & Human system & Level & $N$ & $D$ & Data sparsity (\%) \\
\midrule

Castro-Nallar 2015~\cite{CastroNallar2015} & oral & Genus & 32 & 107 & 60.28 \\
ChngKR 2016~\cite{Chng2016} & skin & Genus & 78 & 211 & 80.66 \\
HMP 2012~\cite{HMP2012} & nasal & Species & 91 & 222 & 92.72 \\
HMP16S v13~\cite{HMP2012} & multi-site 16S & Family & 2909 & 159 & 87.20 \\
American Gut 1~\cite{McDonald2018} & gut citizen science & OTU & 289 & 127 & 30.64 \\
American Gut 2~\cite{McDonald2018} & gut citizen science & OTU & 297 & 138 & 34.48 \\
CosteaPI 2017~\cite{Costea2017} & stool enterotypes & Genus & 279 & 179 & 71.73 \\
MehtaRS 2018~\cite{Mehta2018} & gut longitudinal & Genus & 928 & 183 & 75.72 \\
Diabimmune Karelia~\cite{Vatanen2016} & infant gut & Family & 1584 & 55 & 53.36 \\
ShaoY 2019~\cite{Shao2019} & infant gut & Genus & 1644 & 244 & 93.01 \\
RubelMA 2020~\cite{Rubel2020} & pediatric gut & Species & 175 & 373 & 83.85 \\
KarlssonFH 2013~\cite{Karlsson2013} & stool T2D & Genus & 145 & 158 & 71.06 \\
Schubert CDI~\cite{Schubert2014} & stool CDI & Family & 347 & 79 & 79.35 \\
Lozupone HIV~\cite{Lozupone2013} & gut HIV & Family & 55 & 62 & 65.75 \\
Gevers 2014 IBD~\cite{Gevers2014} & gut IBD & Family & 285 & 72 & 69.85 \\
VilaAV 2018~\cite{Vila2018} & gut IBD/IBS & Genus & 355 & 178 & 77.83 \\
Chan NASH~\cite{Wong2013} & gut NASH & Family & 54 & 24 & 63.12 \\
Baker NASH/obesity~\cite{Zhu2013} & gut NASH/obesity & Family & 64 & 55 & 56.85 \\
Ross obesity~\cite{Ross2015} & gut obesity & Genus & 63 & 123 & 67.05 \\
Zeevi obesity~\cite{Zeevi2015} & gut glycemic response & Family & 820 & 74 & 74.99 \\
Scheperjans PD~\cite{Scheperjans2015} & gut Parkinson's & Family & 148 & 40 & 63.04 \\
Littman RA~\cite{Scher2013} & gut arthritis & Family & 114 & 40 & 67.17 \\
Alkanani T1D~\cite{Alkanani2015} & gut T1D & Family & 112 & 27 & 66.34 \\
YuJ 2015~\cite{Yu2015} & stool CRC & Genus & 128 & 203 & 72.83 \\
Baxter CRC~\cite{Baxter2016} & stool CRC & Family & 490 & 92 & 71.71 \\
Zeller CRC~\cite{Zeller2014} & stool CRC & Family & 129 & 100 & 58.29 \\
Zhao CRC~\cite{Wang2012} & stool CRC & Genus & 102 & 78 & 77.48 \\
MBQC integrated OTUs~\cite{MBQC2017} & stool benchmark & Family & 18364 & 257 & 89.46 \\
\bottomrule
\end{tabular}
\end{table}


\subsection{Real Microbiome Co-occurrence datasets}
\label{subsubsec:network-inference-datasets}
We evaluated network inference on five real microbial interaction benchmarks with experimentally supported edge truth.
Three datasets support broad undirected edge-recovery evaluation: OMM12~\cite{Weiss2022OMM12}, OMM12 keystone 2023~\cite{Weiss2023Keystone}, and PairInteraX~\cite{Zhu2025PairInteraX}.
Two additional datasets support local or directional recovery from experimentally measured pairwise effects: butyrate assembly 2021 and host fitness 2018.
Across these datasets, sample size ranged from $N=109$ to $N=1449$, feature dimension ranged from $D=5$ to $D=96$, and data sparsity ranged from $21.25\%$ to $66.19\%$.
Processed abundance matrices and truth edge sets were harmonized into a common benchmark format so that PLNNetwork and \texttt{GLMNet(Poisson)} could be evaluated under the same scoring rules.
The benchmark datasets are summarized in Table~\ref{tab:network-datasets}.

\begin{table}[t]
\caption{Network-inference datasets used for evaluation against experimental interaction truth.
Data sparsity is the percentage of zero entries in the fitted abundance matrix after count-like preprocessing and removal of zero-sum rows or columns.}
\label{tab:network-datasets}
\centering
\scriptsize
\begin{tabular}{>{\raggedright\arraybackslash}p{0.18\textwidth}>{\raggedright\arraybackslash}p{0.18\textwidth}>{\raggedright\arraybackslash}p{0.24\textwidth}rrr}
\toprule
Dataset & System & Truth type & $N$ & $D$ & Data sparsity (\%) \\
\midrule
OMM12~\cite{Weiss2022OMM12} & defined human gut bacterial community & experimental coculture and dropout interactions & 109 & 12 & 21.25 \\
OMM12 keystone 2023~\cite{Weiss2023Keystone} & defined human gut bacterial community & experimental dropout-derived abundance-ratio effects & 228 & 11 & 31.50 \\
PairInteraX~\cite{Zhu2025PairInteraX} & human gut species panel & experimental pairwise co-culture interaction labels & 169 & 96 & 25.86 \\
Butyrate assembly 2021~\cite{Clark2021Butyrate} & defined human gut community & pair-vs-singleton abundance shifts & 537 & 26 & 66.19 \\
Host fitness 2018~\cite{Gould2018HostFitness} & \textit{Drosophila} gut community & pair-vs-singleton CFU shifts & 1449 & 5 & 54.56 \\
\bottomrule
\end{tabular}
\end{table}



\subsection{Algorithms}
\label{subsec:algorithms}

We used algorithms matched to each task.
All count inputs were log1p-transformed before being passed to any algorithm.
For count prediction, we compared a featureless baseline, \texttt{GLMNet(Poisson)}, and PLN.
For network inference, we compared a diagonal baseline, \texttt{GLMNet(Poisson)} neighbourhood selection, and PLN with graphical lasso.

The featureless baseline (count prediction) predicts each target taxon by the empirical mean of that taxon in the training split, ignoring all cross-taxon information.
The diagonal baseline (network inference) corresponds to an empty graph that predicts no interactions.

\texttt{GLMNet(Poisson)} fits a penalized Poisson regression for each taxon using the abundances of all other taxa as predictors~\cite{Friedman2010}.
For count prediction, we used the lasso penalty with \texttt{cv.glmnet}, modelling each taxon independently with no joint latent covariance.
For network inference, we used nodewise neighbourhood selection~\cite{Meinshausen2006}, retaining coefficients at \texttt{lambda.1se} and defining undirected edges by mutual support after symmetrising the coefficient matrix.

PLN fits a Poisson log-normal model with an intercept and a latent Gaussian layer~\cite{aitchison1989multivariate,Chiquet2021}.
The latent covariance $\hat{\Sigma}$ captures residual dependence among taxa.
For count prediction, predictions for each target taxon are obtained by Gaussian conditioning on the fitted latent distribution.
To stabilize prediction, we apply covariance shrinkage with parameter $\alpha \in \{0.0, 0.1, \ldots, 1.0\}$, where $\alpha$ scales the off-diagonal entries of $\hat{\Sigma}$ as $\hat{\Sigma}_{jk} \leftarrow \alpha\hat{\Sigma}_{jk}$ for $j \neq k$.
For network inference, we apply the graphical lasso~\cite{Friedman2008glasso} to $\hat{\Sigma}$ over 30 log-spaced penalty values from $\rho_{\max}$ to $0.005\,\rho_{\max}$, where $\rho_{\max}$ is the largest off-diagonal entry of $\hat{\Sigma}$ in absolute value.
The sparse precision matrix at the selected penalty defines the estimated interaction graph.

All PLN models were fit with variational inference using the \texttt{PLNmodels} package~\cite{Chiquet2021}.
All penalized Poisson regressions were fit with \texttt{glmnet}~\cite{Friedman2010}.
Graphical lasso fits used the \texttt{glassoFast} package~\cite{Sustik2012glassoFast}.

\subsection{Evaluation protocols}
\label{subsec:evaluation-protocols}

For count prediction, we used leave-one-taxon-out cross-validation (LOTO-CV).
In each outer fold, one taxon is held out as the prediction target and the remaining taxa serve as predictors; this is repeated for every taxon and the scores are averaged.
The outer CV used 3 folds; the inner CV used 3 folds for hyperparameter selection.
The primary metric is held-out Poisson deviance, defined as
\[
  \mathcal{D}(y, \hat\mu) = 2\sum_{i} \left( y_i \log\frac{y_i}{\hat\mu_i} - (y_i - \hat\mu_i) \right),
\]
where $y_i$ are the observed counts and $\hat\mu_i$ are the predicted means.
For \texttt{GLMNet(Poisson)}, the regularization parameter $\lambda$ was selected by inner-fold Poisson deviance using \texttt{cv.glmnet}.
For PLN, the shrinkage parameter $\alpha$ was selected by 3-fold inner cross-validation on Poisson deviance.
The featureless baseline serves as a lower bound: any algorithm with higher deviance than the featureless baseline provides no useful predictive signal.

For network inference, estimated interaction graphs were compared against experimentally validated microbial interactions from five ground-truth datasets: PairInteraX~\cite{Zhu2025PairInteraX}, OMM12~\cite{Weiss2022OMM12}, OMM12 keystone 2023~\cite{Weiss2023Keystone}, butyrate assembly 2021~\cite{Clark2021Butyrate}, and host fitness 2018~\cite{Gould2018HostFitness}.
Each dataset provides experimentally derived pairwise interaction labels obtained from co-culture, dropout, or combinatorial abundance experiments.
The primary metrics are precision, recall, and F1 score on edge recovery, computed by matching predicted edges against the ground-truth interaction set.
For PLN with graphical lasso, the penalty $\rho$ was selected by EBIC on the training data.
For \texttt{GLMNet(Poisson)} neighbourhood selection, $\lambda$ was fixed at \texttt{lambda.1se} from \texttt{cv.glmnet}.
The diagonal baseline, which predicts an empty graph, serves as the trivial reference point.

\subsection{Computational setup}
\label{subsec:computational-setup}

All experiments were run on Northern Arizona University's Monsoon high-performance computing cluster~\cite{Buechler2025MonsoonIntro}.
Prediction and network benchmark jobs were executed as single-core Slurm jobs.
PLN-based learners were additionally run with one Torch thread to avoid conflating algorithmic differences with thread-level speedups.

Runtime and memory scaling were profiled using the \texttt{atime} R package on the Vogtmann et al.\ stool species dataset~\cite{Vogtmann2016CRC}.
For the saved profiling run, the number of taxa $p$ was varied from 10 to 1000 at a fixed sample size of $N=110$, and each configuration was run three times.
We report median wall time and peak memory across those repeated measurements.


\section{Results}
\subsection{Variation Across Datasets in Relative Algorithm Performance}

\begin{figure}[t]
\centering
\begin{tikzpicture}
\begin{axis}[
  width=0.62\textwidth,
  height=0.38\textwidth,
  xlabel={$N/D$ ratio (log scale)},
  ylabel={Mean absolute correlation (MAC)},
  xmode=log,
  xmin=0.05, xmax=100,
  ymin=0, ymax=0.25,
  ytick={0, 0.05, 0.10, 0.15, 0.20, 0.25},
  yticklabel style={/pgf/number format/fixed, /pgf/number format/fixed zerofill,
                    /pgf/number format/precision=2, font=\footnotesize},
  ymajorgrids=true,
  xmajorgrids=true,
  grid style={dotted, gray!40},
  tick label style={font=\footnotesize},
  label style={font=\footnotesize},
  legend style={font=\footnotesize, at={(0.97,0.97)}, anchor=north east},
  legend cell align=left,
]
\addplot[only marks, mark=*, mark size=2.2pt, blue!75!black]
  coordinates {
    (0.30,0.167)(0.37,0.092)(0.41,0.118)(0.42,0.117)(0.47,0.056)
    (0.51,0.130)(0.63,0.063)(0.89,0.163)(0.92,0.067)
    (1.16,0.156)(1.31,0.091)(1.76,0.056)(2.13,0.237)(2.25,0.227)
    (2.85,0.112)(3.70,0.079)(3.96,0.085)(4.15,0.231)(4.39,0.074)
    (5.33,0.086)
  };
\addlegendentry{PLN wins ($n=21$)}
\addplot[only marks, mark=square*, mark size=2.2pt, red!75!black]
  coordinates {
    (1.40,0.173)(1.56,0.055)(1.99,0.049)(2.28,0.182)
    (6.74,0.033)(11.08,0.110)(18.30,0.058)
    (28.80,0.127)(77.74,0.041)
  };
\addlegendentry{\texttt{GLMNet} wins ($n=9$)}
\end{axis}
\end{tikzpicture}
\caption{Dataset-level winner across 30 real microbiome datasets as a function
of the sample-to-taxon ratio $N/D$ and mean absolute correlation (MAC) among taxa.
PLN wins predominantly in the upper-left region (low $N/D$, high MAC);
\texttt{GLMNet(Poisson)} wins predominantly in the lower-right region (high $N/D$, low MAC).}
\label{fig:mac-nd-2d}
\end{figure}

Neither PLN nor GLMNet(Poisson) uniformly dominates across the 30 benchmark datasets
(Figure~\ref{fig:mac-nd-2d}).
The sample-to-taxon ratio $N/D$ is the primary determinant of the outcome, yielding an AUC of 0.83 for predicting the dataset-level winner on its own.
The mean absolute correlation (MAC) among taxa, shown on the $y$-axis of Figure~\ref{fig:mac-nd-2d}, provides an intuitive account of the pattern: PLN tends to win in the upper-left region where $N/D$ is small and inter-taxon correlations are strong, while GLMNet(Poisson) wins in the lower-right where $N/D$ is large and correlations are weak.
Statistically, however, overdispersion (the mean variance-to-mean count ratio across taxa) is a stronger secondary predictor than MAC, raising AUC from 0.83 to 0.93 when combined with $N/D$, compared to 0.84 for MAC.
This is consistent with PLN's design: the latent Gaussian layer explicitly absorbs extra-Poisson variance, so its advantage grows with the degree to which counts deviate from the Poisson assumption.


\begin{figure}[t]
  \centering
  \includegraphics[width=\textwidth]{count_prediction_case_study_benchmark}
  \caption{Representative datasets illustrating both directions of the PLN
  vs.\ GLMNet(Poisson) comparison.
  Top: datasets where PLN achieves lower mean Poisson deviance.
  Bottom: datasets where GLMNet(Poisson) achieves lower mean Poisson deviance.
  Each panel reports mean Poisson deviance (mean $\pm$ SE across three outer
  cross-validation folds) for PLN, GLMNet(Poisson), and the featureless baseline.
  The annotated difference and $p$-value are from the paired comparison of PLN
  and GLMNet(Poisson) across folds.}
  \label{fig:count_pred_case_study}
\end{figure}

Figure~\ref{fig:count_pred_case_study} illustrates both directions with representative datasets.
In the PLN-winning cases, Lozupone HIV family ($N/D = 0.89$), Zeller CRC genus ($N/D = 0.50$), Schubert CDI family ($N/D = 4.39$), and amgut2 ($N/D = 2.13$) all show PLN achieving substantially lower deviance than GLMNet(Poisson), with differences ranging from 0.08 to 0.49 deviance units and $p < 0.001$ in each case.
In the GLMNet(Poisson)-winning cases, CosteaPI 2017 stool ($N/D = 1.56$, low overdispersion) and the two large-sample datasets MBQC integrated ($N=18{,}270$) and Diabimmune Karelia ($N=1{,}584$) all show GLMNet(Poisson) achieving lower deviance, with PLN offering no improvement despite the available sample sizes.
In all eight representative datasets, both algorithms improve upon the featureless baseline.


\textbf{Implications for algorithm selection.}
The choice between PLN and GLMNet(Poisson) depends on properties of the data that are not always known in advance, including the strength and structure of inter-taxon covariance, the sparsity pattern of the count matrix, and the signal-to-noise ratio of the prediction task, all of which interact with $N$, $D$, and their ratio.
A practical guideline emerges from the results: when $N/D$ is low or overdispersion is high, PLN is the stronger choice; when $N/D$ is large or counts are near-Poisson, GLMNet(Poisson) is competitive and considerably faster.
Practitioners should treat these as starting heuristics and empirically validate algorithm choice on their specific datasets.


\subsection{Network-inference results}
\label{subsec:network-inference-results}


\begin{figure}[t]
  \centering
  \includegraphics[width=\textwidth]{interaction_truth_benchmarks_f1_2026-03-28}
  \caption{Edge-recovery F1 on five experimental ground-truth datasets for
  PLNNetwork and GLMNet(Poisson).
  Datasets are grouped by ground-truth type: broad edge recovery (OMM12,
  OMM12 keystone 2023, PairInteraX) and local/directional (butyrate assembly 2021,
  host fitness 2018).}
  \label{fig:network_f1}
\end{figure}

\begin{figure}[t]
  \centering
  \includegraphics[width=\textwidth]{interaction_network_example_omm12_2026-03-28}
  \caption{Ground-truth and inferred networks for the OMM12 dataset.
  Each panel shows the experimental ground truth (left), the PLNNetwork prediction
  (centre), and the GLMNet(Poisson) prediction (right).
  Edge colour indicates interaction sign (blue: positive, red: negative, grey: mixed).}
  \label{fig:network_case_studies}
\end{figure}

PLNNetwork outperforms GLMNet(Poisson) on four of five ground-truth datasets,
but the margin and direction depend on the ground-truth type (Figure~\ref{fig:network_f1}).
On the three broad edge recovery datasets, PLNNetwork achieves substantially
higher F1: OMM12 ($\text{F1} = 0.86$ vs.\ $0.29$), OMM12 keystone 2023
($0.62$ vs.\ $0.21$), and PairInteraX ($0.23$ vs.\ $0.02$).
The OMM12 case study (Figure~\ref{fig:network_case_studies}) illustrates why:
the experimental ground truth is a dense, predominantly negative interaction
network that PLNNetwork recovers well, while GLMNet(Poisson) infers a sparse
network with predominantly positive edges.
On the local/directional datasets, the pattern reverses for host fitness 2018:
GLMNet(Poisson) achieves $\text{F1} = 0.95$ vs.\ PLNNetwork's $0.80$,
recovering all nine experimental interactions with a single false positive
.
PLNNetwork correctly identifies six of the nine interactions but misses three
edges involving taxon AO.
On butyrate assembly 2021, PLNNetwork retains a small advantage
($\text{F1} = 0.36$ vs.\ $0.27$).

\subsection{Computational scaling}
\label{subsec:computational-scaling}

\begin{figure}[t]
  \centering
  \includegraphics[width=\textwidth]{vogtmann_stool_species_pln_glmnet_p_scaling}
  \caption{Median wall time (left) and peak memory (right) as a function of the
  number of taxa $D$, profiled on the Vogtmann et al.\ stool species dataset
  ($N = 110$) using the \texttt{atime} R package.
  Each point is the median of three runs; $D$ ranges from 10 to 1000.
  Both axes are on a log scale.}
  \label{fig:atime}
\end{figure}

PLN is substantially more expensive than GLMNet(Poisson) in both wall time and
peak memory at all values of $D$ (Figure~\ref{fig:atime}).
Wall time for PLN grows approximately as $O(D^2)$ to $O(D^3)$, reaching 76\,s
at $D = 1000$, while GLMNet(Poisson) scales approximately as $O(D)$, reaching
0.28\,s, a roughly 280-fold difference at the largest setting.
Peak memory shows a similar gap: PLN uses 392\,MB at $D = 1000$, compared to
79\,MB for GLMNet(Poisson).
Despite the steeper scaling, PLN remains feasible for the dataset sizes
typical of microbiome benchmarks: wall time stays below 80\,s and peak memory
below 400\,MB across the full range of $D$ tested.


\section{Discussion}
\label{sec:discussion}

The dominance of $N/D$ as a predictor of relative performance is consistent with classical estimation theory: PLN's latent covariance matrix has $O(D^2)$ free parameters, so reliable estimation requires $N$ to grow proportionally with $D$.
When this condition is not met, the implicit regularisation of GLMNet(Poisson)'s independent-taxon lasso regressions is better adapted to the available signal.
The secondary role of overdispersion is equally direct: PLN's latent Gaussian layer is specifically designed to absorb extra-Poisson variance, so when count distributions are near-Poisson, that layer adds parameters without improving fit.
Taken together, $N/D$ and overdispersion reach an AUC of 0.93 for predicting the dataset-level winner, suggesting that practitioners can form a reliable prior expectation of relative performance from properties of the count matrix alone, without fitting either model.
The result also offers an empirical rationalisation of why mean absolute correlation (MAC) among taxa, while intuitive, adds less predictive power than overdispersion: MAC captures the structure of the covariance that PLN could exploit, whereas overdispersion captures whether that exploitation is mechanistically warranted.

The network-inference results reflect a structural difference between the two
approaches.
PLN recovers an undirected graphical model by applying graphical lasso to the
estimated latent covariance; this is well-suited to broad, symmetric interaction
networks such as OMM12, where the ground truth encodes pairwise competition and
facilitation among all member species.
GLMNet(Poisson) fits a separate penalised regression for each taxon, following
the neighbourhood selection approach of \citet{Meinshausen2006}, and is better
aligned with local or directional interactions where the effect of one species
on another is not necessarily reciprocal.
The host fitness ground truth, which records the directional effect of individual
taxa on host reproductive fitness, exemplifies this setting.

Several limitations constrain the generalisability of these findings.
The benchmark compares only PLN and GLMNet(Poisson) and does not evaluate other
methods for microbial network inference, such as SPIEC-EASI~\citep{Kurtz2015SPIECEASI}
and SparCC~\citep{Friedman2012SparCC}, which are designed to account for the
compositional structure of relative-abundance data.
The network-inference evaluation is restricted to five experimental ground-truth
datasets; results may not generalise to communities with different sizes or
interaction densities.
Both algorithms receive the same log1p-transformed input, a preprocessing choice
that may not be optimal for PLN, which is designed for raw count data.
The computational profiling is limited to a single dataset at fixed $N = 110$
on a single core; results may differ at larger sample sizes or with multi-threaded
execution.

\section{Conclusion}
\label{sec:conclusion}

This study provides the first systematic empirical map of when PLN outperforms penalized Poisson regression on real microbiome data, evaluated under a unified held-out protocol across 34 datasets and five experimentally validated network ground truths.
Neither model uniformly dominates.
For count prediction, $N/D$ is the primary signal for algorithm selection and overdispersion is the strongest secondary discriminant: PLN is the stronger choice when samples are few relative to taxa and counts are substantially overdispersed, while GLMNet(Poisson) is competitive and considerably faster when $N/D$ is large.
For network inference, PLNNetwork better recovers broad, undirected interaction graphs, while GLMNet(Poisson) nodewise regression is better suited to local or directional effects.
Despite a two-to-three order-of-magnitude wall-time difference at large $D$, PLN remains within practical limits for current microbiome dataset sizes.
These findings give practitioners actionable, data-driven criteria for choosing between the two model classes without requiring both to be fitted.

\paragraph{Key findings.}
The benchmark yields four main observations.
First, $N/D$ is the primary predictor of the dataset-level winner: PLN dominates when samples are few relative to taxa, and GLMNet(Poisson) dominates when the ratio is large.
Second, overdispersion is the strongest secondary signal, outperforming mean absolute correlation when combined with $N/D$ (AUC 0.93 vs.\ 0.84), consistent with PLN’s latent Gaussian layer being most beneficial when count distributions substantially deviate from the Poisson assumption.
Third, for network inference, the structural difference between the two approaches matters: PLNNetwork recovers broad, undirected interaction graphs more reliably, while GLMNet(Poisson) nodewise regression is better suited to local or directional effects.
Fourth, PLN’s computational cost is two to three orders of magnitude higher than GLMNet(Poisson) at large $D$, yet remains within practical limits for the dataset sizes typical of current microbiome studies.

\paragraph{Future directions.}
Several directions follow from the limitations of this study.
First, the benchmark draws from a finite set of public datasets; broader coverage of environments such as soil, marine, and host-associated systems, as well as diverse sequencing protocols, would strengthen the generalisability of the conclusions.
Second, all algorithms received the same \texttt{log1p}-transformed input~\cite{doi:10.1128/mSystems.00016-19}; different normalisation or offset strategies may shift relative performance, and evaluating PLN on raw counts (the input it is designed for) is an important next step.
Third, the comparison is restricted to algorithms in the Poisson family; extending the protocol to negative-binomial and zero-inflated models would give a more complete picture of the count-modelling landscape~\cite{10.1038/s41598-024-76513-8,10.1101/538579,10.1093/bioadv/vbae167,Champion2024OneNet}.
Fourth, lightweight diagnostics that estimate recoverable latent dependence (for example, through stability of conditional correlations across folds) could help practitioners decide when PLN is likely to provide benefits before incurring its full computational cost.
Finally, the computational experiments were conducted on Northern Arizona University’s Monsoon supercomputer~\cite{Buechler2025MonsoonIntro}; profiling PLN under multi-threaded execution and at larger sample sizes would give a more complete picture of its practical scalability.

\bibliographystyle{plainnat}
\bibliography{references}
\end{document}